\section{Parameter splittings and universal analytic reflection}\label{app:reflection}\label{fam:sec:tensors}
A normed space $P$ is polarizable in degree $n$ (condition $(\Gamma^n_P)$, Definition~\ref{ob:def:sequences}) if a bounded
projection $R_n$ of $\widehat\otimes^n_\pi P$ onto $\widehat\Gamma^n_\pi P$ exists; for Banach $P$ over complete $K$ this is
reflection in degree $n$ (Proposition~\ref{fam:prop:tensor-homogeneous}). Analytic descent (Section~\ref{sec:descent}) of every
Banach-valued analytic map through every closed inclusion is stronger; we characterize it and separate the two by an example.

Throughout this appendix, $K$ is complete and nontrivially normed, and $P$ is Banach.
For $n\ge1$, let $\widehat\otimes^n_\pi P$ be the completed projective tensor power as in
Section~\ref{ob:ssec:powers}; it represents bounded $n$\nobreakdash-linear maps from $P^n$ to Banach spaces.
Recall from Definition~\ref{ob:def:sequences} the closed subspace
\[
 \widehat\Gamma^n_\pi P=\overline{\spn\{x^{\otimes n}:x\in P\}}\subseteq \widehat\otimes^n_\pi P.
\]

\begin{proposition}[homogeneous reflection]\label{fam:prop:tensor-homogeneous}
For a fixed $n\ge1$, the following conditions are equivalent.
\begin{enumerate}
\item For every Banach space $Z$, closed subspace $W\subseteq Z$, and bounded $n$\nobreakdash-linear map
$B\colon P^n\to Z$ with diagonal values in $W$, there is a bounded $W$\nobreakdash-valued
$n$\nobreakdash-linear map with the same diagonal.
\item The subspace $\widehat\Gamma^n_\pi P$ is the range of a bounded projection in $\widehat\otimes^n_\pi P$.
\end{enumerate}
\end{proposition}

\begin{proof}
Suppose first that a projection $R_n$ exists. Linearize $B$ to a bounded map
$\widetilde B\colon \widehat\otimes^n_\pi P\to Z$. The diagonal hypothesis and closedness of $W$ imply
$\widetilde B(\widehat\Gamma^n_\pi P)\subseteq W$.
Composing $\widetilde B|_{\widehat\Gamma^n_\pi P}$ with $R_n$ and the canonical tensor map gives the required
multilinear representative, of norm at most $\norm B\norm{R_n}$.
Conversely, take $Z=\widehat\otimes^n_\pi P$, $W=\widehat\Gamma^n_\pi P$, and $B$ equal to the canonical tensor map.
Linearizing the representative supplied by the first condition gives a bounded map into $\widehat\Gamma^n_\pi P$
fixing every pure power. It fixes their linear span and, by continuity, their closed span, so it is a projection.
\end{proof}

\begin{definition}\label{fam:def:universal-reflection}
The Banach parameter space $P$ has \emph{universal analytic reflection} if, for every Banach space $Z$, closed
subspace $W\subseteq Z$, every open $U\subseteq P$, and every map $f\colon U\to W$, analyticity of
the composite $U\to W\hookrightarrow Z$ implies analyticity of $f$.
\end{definition}

\begin{theorem}[analytic tensor criterion]\label{fam:thm:tensor-analytic}
A Banach space $P$ has universal analytic reflection if and only if there are projections
$R_n\colon \widehat\otimes^n_\pi P\to\widehat\Gamma^n_\pi P$ for all $n\ge1$ such that
\begin{equation}\label{fam:eq:exponential-projections}
 \limsup_{n\to\infty}\norm{R_n}^{1/n}<\infty.
\end{equation}
\end{theorem}

\begin{proof}
Assume first that the projections exist. Expand an ambient analytic map with values in $W$ at a point.
The scalar-line uniqueness argument in Theorem~\ref{thm:coefficient-descent} puts every homogeneous diagonal in $W$.
The construction in Proposition~\ref{fam:prop:tensor-homogeneous} replaces its coefficient $p_n$ by a
$W$\nobreakdash-valued coefficient of norm at most $\norm{p_n}\norm{R_n}$ and with the same diagonal.
The exponential bound preserves a positive radius of convergence. The degree-zero coefficient is
already in $W$, and the resulting expansion represents the given map.

Conversely, form the Banach spaces with supremum norms
\[
 Z=c_0\bigl(\{\widehat\otimes^n_\pi P\}_{n\ge1}\bigr),\qquad
 W=c_0\bigl(\{\widehat\Gamma^n_\pi P\}_{n\ge1}\bigr)\subseteq Z.
\]
Here $c_0$ denotes families tending to zero in their respective Banach coordinate spaces; $W$ is closed.
On the open unit ball of $P$, define
\[
 f(x)=(x^{\otimes n})_{n\ge1}.
\]
Since $\norm{x^{\otimes n}}\le\norm x^n$, this vector lies in $W$.
As a $Z$\nobreakdash-valued map it is analytic: its degree-$n$ coefficient is the canonical tensor map
placed in coordinate $n$, and has norm at most one. This gives an expansion at zero. To re-expand
at $a$ with $\norm a<1$, choose $r>0$ with $\norm a+r<1$. Expanding each tensor power in $a+h$
gives translated coefficients, and the sum of all term bounds is at most
$\sum_{n\ge1}(\norm a+r)^n<\infty$. Completeness of $Z$ supplies the coefficients and the
rearranged sum, proving analyticity throughout the unit ball. Universal reflection gives a $W$\nobreakdash-valued
expansion $\sum_mq_m(x,\ldots,x)$ near zero with $\norm{q_m}\le C\rho^{-m}$ for some $C,\rho>0$.
Compose this expansion with the projection $\pi_n\colon W\to\widehat\Gamma^n_\pi P$ and compare on scalar lines
with $x\mapsto x^{\otimes n}$. Uniqueness of scalar power-series coefficients gives
\[
 \pi_nq_n(x,\ldots,x)=x^{\otimes n}\qquad(x\in P).
\]
Linearizing $\pi_nq_n$ therefore gives a projection $R_n$ onto $\widehat\Gamma^n_\pi P$ with
$\norm{R_n}\le\norm{q_n}\le C\rho^{-n}$, proving~\eqref{fam:eq:exponential-projections}.
\end{proof}

\subsection{Projection bounds and the special role of fixed degree}
For $P=K^d$ with the maximum norm, the coordinate construction in
Theorem~\ref{thm:finite-parameters}, applied to the canonical tensor map, gives
$\norm{R_n}\le d^n$ for $d\ge1$. A bounded linear change of coordinates introduces only a fixed
distortion factor to the $n$\nobreakdash-th power, so it does not change the exponential-growth condition.

For $P=\ell^1(I,K)$, the proof of Proposition~\ref{spl:prop:gamma-l1}(1) identifies $\widehat\otimes^n_\pi P$
isometrically with $\ell^1(I^n,K)$, with unit vectors $e_w$ for $w\in I^n$, puts the sum of the $e_w$ over each
permutation orbit in $I^n$ into $\widehat\Gamma^n_\pi P$, and gives a projection onto $\widehat\Gamma^n_\pi P$ of norm
at most $n!$. In these coordinates every element of $\widehat\Gamma^n_\pi P$ is constant on each permutation orbit,
since every pure power is.
Over a nonarchimedean field and for infinite $I$, the bound $n!$ is optimal even among projections that
do not preserve the individual orbit blocks.

\begin{proposition}\label{fam:prop:l1-optimal}
If $K$ is nonarchimedean and complete, and $I$ is infinite, then
\[
 \inf\{\norm R:R\colon \widehat\otimes^n_\pi\ell^1(I,K)\to\widehat\Gamma^n_\pi\ell^1(I,K)
                  \text{ is a bounded projection}\}=n!.
\]
Consequently $\ell^1(I,K)$ does not have universal analytic reflection.
\end{proposition}

\begin{proof}
Proposition~\ref{spl:prop:gamma-l1}(1) gives the upper bound. For the lower bound, take an orbit $\mathcal O$
of a tuple with $n$ distinct indices; it has $n!$ elements. For each $w\in\mathcal O$, let $\lambda_w$
be the common coefficient on $\mathcal O$ of $R(e_w)$.
The orbit sum $\sum_{w\in\mathcal O}e_w$ belongs to $\widehat\Gamma^n_\pi P$ and is fixed by $R$, hence
\[
 \sum_{w\in\mathcal O}\lambda_w=1.
\]
The nonarchimedean inequality implies that some $\abs{\lambda_w}\ge1$.
For this $w$, the ordinary sum norm satisfies
$\norm{R(e_w)}_1\ge n!\abs{\lambda_w}\ge n!$.
Thus $\norm R\ge n!$. Finally, $(n!)^{1/n}$ is unbounded, so
Theorem~\ref{fam:thm:tensor-analytic} excludes universal analytic reflection.
\end{proof}

\begin{corollary}[entire polynomial approximation of a nonanalytic smooth map]
\label{fam:cor:no-common-radius}
Let $K$ be a complete nonarchimedean nontrivially normed field and put
$P=\ell^1(\N,K)$, with the ordinary sum norm. There are Banach spaces
$W\subseteq Z$, with $W$ closed, and a map $f$ on the open unit ball of $P$
such that $f$ is analytic into $Z$ and $C^\infty$ into $W$, but is not analytic
into $W$ at zero. Moreover, $f$ is a uniform limit on every smaller ball of
entire $W$\nobreakdash-valued polynomials.
\end{corollary}

\begin{proof}
Take
\[
 Z=c_0\bigl(\{\widehat\otimes^n_\pi P\}_{n\ge1}\bigr),\qquad
 W=c_0\bigl(\{\widehat\Gamma^n_\pi P\}_{n\ge1}\bigr),\qquad
 f(x)=(x^{\otimes n})_{n\ge1}.
\]
The spaces are Banach and $W$ is closed in $Z$. As in the proof of
Theorem~\ref{fam:thm:tensor-analytic}, the ambient expansion has coefficient
norms at most one and defines an analytic map on the open unit ball.
Smooth reflection through the closed isometric inclusion makes $f$
$C^\infty$ as a map into $W$.

Suppose a $W$\nobreakdash-valued expansion at zero has coefficients $q_n$.
Uniqueness on scalar lines gives
\[
 \pi_nq_n(x,\ldots,x)=x^{\otimes n},
\]
where $\pi_n\colon W\to\widehat\Gamma^n_\pi P$ is the coordinate projection.
Linearizing $\pi_nq_n$ gives a projection
$\widehat\otimes^n_\pi P\to\widehat\Gamma^n_\pi P$ of norm at most $\norm{q_n}$.
Proposition~\ref{fam:prop:l1-optimal} therefore gives
$\norm{q_n}\ge n!$. Since $(n!)^{1/n}$ is unbounded, these coefficients
give a formal multilinear series of radius zero, contradicting the assumed analytic expansion.

For each $n$, choose the sorted-orbit projection
$R_n\colon \widehat\otimes^n_\pi P\to\widehat\Gamma^n_\pi P$. Its composition with the canonical
tensor map is a bounded multilinear representative of $x\mapsto x^{\otimes n}$
with values in $\widehat\Gamma^n_\pi P$. Consequently the truncation
\[
 f_N(x)=(x,x^{\otimes2},\ldots,x^{\otimes N},0,\ldots)
\]
is an entire $W$\nobreakdash-valued polynomial. For $0<r<1$ and $\norm x\le r$,
\[
 \norm{f(x)-f_N(x)}_W\le r^{N+1},
\]
which proves the asserted uniform convergence.
\end{proof}


There is no conflict with Theorem~\ref{thm:summable-parameters}(2): the ambient representation there has one fixed
degree $d$, so factoring a general analytic parameter map through a free $\ell^1$ space incurs a single factorial
bound $d!$, rather than a sequence of factorial bounds across Taylor degrees.
Over $\mathbb R$ and $\mathbb C$ every Banach space has universal analytic reflection, by Theorem~\ref{fam:thm:tensor-analytic}: the symmetrization $\frac1{n!}\sum_{\sigma\in S_n}\sigma$ has norm at most one on $\widehat\otimes^n_\pi P$ and, by polarization, is a projection onto $\widehat\Gamma^n_\pi P$.
